\section{27.02.2026}{Reprezentacje unitarne}

\begin{definition}{}{}
  ${}^*$-podalgebra $A\subseteq \B(\Hh)$ nazywa się \buff{algebrą von Neumanna}, jeśli
  \begin{enumerate}
    \item $A$ jest domknięte w WO
    \item $A=A''$
  \end{enumerate}
  gdzie $X':=\{T\in\B(\Hh)\;:\;\forall\;s\in X\;Ts=sT\}$ to \acc{komutant} zbioru $X\subseteq \B(\Hh)$.
\end{definition}

\begin{fact}{}{}
  Jeśli $X=X^*$ to $X'$ jest algebrą von Neumanna.
\end{fact}

\begin{proof}
  Wychodzimy od stwierdzenia, że
  $$X\subseteq X'',$$
  które widać od razu rozpisując definicje
  $$X'=\{T\;:\;\forall\; s\in X\;Ts=sT\}$$
  $$X''=\{T\;:\;\forall\; s'\in X\;Ts'=s'T\}$$

  Stawiając $X'$ pod $X$ dostajemy $X'\subseteq X'''=(X')''$.

  Z drugiej strony, nakładanie $'$ odwraca zawierania (łatwiej komutować z każdym elementem mniejszego zbioru), więc $X'\supseteq X'''=(X')''$.
\end{proof}

\begin{fact}{}{}\color{red}
  Jeśli $T=T^*$ to $T$ ma rozkład spektralny
  $$T=\int t dE_t$$
\end{fact}

\begin{lemma}{Schura}{}
  \begin{enumerate}
    \item Jeśli $\pi$ jest reprezentacją nieprzywiedlnia to $Hom(\pi, \pi)=\C$.
    \item Jeśli $\pi$ oraz $\pi'$ są nieprzywiedlnie, to $\dim Hom(\pi, \pi')\leq 1$.
  \end{enumerate}
\end{lemma}

Punkt 2. oznacza w szczególności, że albo dwie reprezentacje nieprzywiedlnie są izomorficzne, albo wogóle nie mają przeplataczy.

\begin{proof}
  \color{red}tutaj trzeba wstawkę z twierdzeniem spektralnym
\end{proof}

Wynika z tego, że przeplatacze reprezentacji nad $\C$ nie mają jądra, czyli tworza algebrę z dzieleniem. Nad $\R$ jest trochę trudniej i przeplatacze przychodzą w trzech smakach
\begin{itemize}
  \item krotność identyczności,
  \item obroty (struktura zespolona),
  \item struktura kwaternionowa.
\end{itemize}

\begin{conclusion}{}{}
  Nieprzywiedlnie zespolone reprezentacje grup abelowych są jednowymiarowe.
\end{conclusion}

\begin{proof}
  Niech $\pi\subseteq \pi'$ będzie podreprezentacją. Wtedy dla każdego $g$ 
\end{proof}










